\documentclass[11pt]{article}

\usepackage[margin=1in]{geometry}
\usepackage{enumitem}
\usepackage{hyperref}
\usepackage{xcolor}
\usepackage{titlesec}
\usepackage{listings}
\usepackage[T1]{fontenc}
\usepackage{lmodern}
\usepackage{parskip}
\usepackage{fancyhdr}
\usepackage{microtype}
\usepackage{graphicx}
\usepackage{longtable}
\usepackage{booktabs}
\usepackage{array}

% ── Colors ──────────────────────────────────────────────
\definecolor{accent}{HTML}{1A3C6E}
\definecolor{codebg}{HTML}{F5F5F5}
\definecolor{codeframe}{HTML}{CCCCCC}
\definecolor{sqlkw}{HTML}{0550AE}
\definecolor{sqlcomment}{HTML}{6A737D}
\definecolor{routecolor}{HTML}{0550AE}

% ── Hyperlinks ──────────────────────────────────────────
\hypersetup{
    colorlinks=true,
    linkcolor=accent,
    urlcolor=accent,
}

% ── Section styling ─────────────────────────────────────
\titleformat{\section}
    {\Large\bfseries\color{accent}}
    {\thesection.}{0.5em}{}
    [\vspace{-0.4em}\textcolor{accent}{\rule{\textwidth}{0.8pt}}]
\titlespacing{\section}{0pt}{1.5em}{0.8em}

\titleformat{\subsection}
    {\large\bfseries\color{accent!80!black}}
    {\thesubsection.}{0.5em}{}
\titlespacing{\subsection}{0pt}{1em}{0.4em}

\titleformat{\subsubsection}
    {\normalsize\bfseries\color{accent!60!black}}
    {\thesubsubsection.}{0.5em}{}
\titlespacing{\subsubsection}{0pt}{0.8em}{0.3em}

% ── Header / Footer ────────────────────────────────────
\pagestyle{fancy}
\fancyhf{}
\renewcommand{\headrulewidth}{0.4pt}
\fancyhead[L]{\small\textcolor{accent}{Milestone 4: API Specification}}
\fancyhead[R]{\small\textcolor{accent}{Gavini, Lee, and Ni}}
\fancyfoot[C]{\small\thepage}

% ── Code listing style ──────────────────────────────────
\lstdefinestyle{sql}{
    language=SQL,
    basicstyle=\ttfamily\small,
    breaklines=true,
    frame=single,
    rulecolor=\color{codeframe},
    backgroundcolor=\color{codebg},
    keywordstyle=\color{sqlkw}\bfseries,
    commentstyle=\color{sqlcomment}\itshape,
    showstringspaces=false,
    tabsize=4,
    xleftmargin=1em,
    framexleftmargin=0.5em,
    aboveskip=1em,
    belowskip=1em,
    morekeywords={SERIAL,BIGSERIAL,BIGINT,TIMESTAMPTZ,TEXT,BOOLEAN,DOUBLE,PRECISION,CASCADE,ILIKE,COALESCE,INTERVAL},
}

% ── Route command ───────────────────────────────────────
\newcommand{\route}[2]{\texttt{\textcolor{routecolor}{\textbf{#1}}~#2}}

% ── List styling ────────────────────────────────────────
\setlist[itemize]{leftmargin=1.5em, itemsep=0.3em}
\setlist[enumerate]{leftmargin=1.5em, itemsep=0.4em}

\begin{document}

% ── Title page (no header/footer) ──────────────────────
\thispagestyle{empty}
\begin{center}
    \vspace*{3cm}
    {\LARGE\bfseries\color{accent} Milestone 4: API Specification}\\[0.5em]
    {\Large Stock News Trader}\\[0.8em]
    \textcolor{accent}{\rule{0.6\textwidth}{1pt}}\\[1em]
    {\large Kartheek Gavini, Philip Lee, and Na Ni}
\end{center}

\vfill

%═══════════════════════════════════════════════════════════
\newpage
\section{Routes Overview}
%═══════════════════════════════════════════════════════════

\begin{center}
\small
\begin{longtable}{c l l p{5.5cm} l}
\toprule
\textbf{\#} & \textbf{Method} & \textbf{Path} & \textbf{Description} & \textbf{Type} \\
\midrule
\endfirsthead
\toprule
\textbf{\#} & \textbf{Method} & \textbf{Path} & \textbf{Description} & \textbf{Type} \\
\midrule
\endhead
\bottomrule
\endfoot
1  & GET  & \texttt{/api/companies}                  & List all companies with sector/industry              & Simple \\
2  & GET  & \texttt{/api/companies/search}            & Search companies by ticker or name                   & Simple \\
3  & GET  & \texttt{/api/companies/:ticker}            & Company profile with latest price                    & Simple \\
4  & GET  & \texttt{/api/companies/:ticker/prices}     & Historical OHLCV price data for charting             & Simple \\
5  & GET  & \texttt{/api/stocks/top-gainers}           & Top 10 gainers on a given date (CTE + window)        & \textbf{Complex} \\
6  & GET  & \texttt{/api/sectors/:sector/companies}    & Companies filtered by sector                         & Simple \\
7  & GET  & \texttt{/api/companies/:ticker/news}       & Latest news articles for a company                   & Simple \\
8  & GET  & \texttt{/api/news/trending}                & Most-mentioned companies in recent news              & \textbf{Complex} \\
9  & GET  & \texttt{/api/sectors/returns}              & Average return by sector on a given date (CTE + agg) & \textbf{Complex} \\
10 & GET  & \texttt{/api/sectors}                      & List all sectors                                     & Simple \\
A1 & POST & \texttt{/api/users}                        & Create a new user account                            & Auxiliary \\
A2 & POST & \texttt{/api/users/login}                  & Authenticate a user                                  & Auxiliary \\
\end{longtable}
\end{center}

%═══════════════════════════════════════════════════════════
\newpage
\section{Route Definitions}
%═══════════════════════════════════════════════════════════

%───────────────────────────────────────────────────────────
\subsection{Route 1: List All Companies}
%───────────────────────────────────────────────────────────

\route{GET}{/api/companies}

\medskip\noindent
\textbf{Description:} Listing all companies with their sector and industry names. Maps to Milestone~3 Query~1.

\subsubsection*{Request Parameters}

None.

\subsubsection*{SQL Query}
\begin{lstlisting}[style=sql,basicstyle=\ttfamily\scriptsize]
SELECT c.company_id, c.ticker, c.company_name,
       i.industry_name, s.sector_name
FROM company c
JOIN industry i ON c.industry_id = i.industry_id
JOIN sector s   ON i.sector_id   = s.sector_id
ORDER BY c.ticker
\end{lstlisting}

\subsubsection*{Response Schema}

\begin{center}
\small
\begin{tabular}{l l p{7cm}}
\toprule
\textbf{Field} & \textbf{Data Type} & \textbf{Description} \\
\midrule
\texttt{company\_id}    & integer & Internal company identifier \\
\texttt{ticker}         & string  & Stock ticker symbol \\
\texttt{company\_name}  & string  & Full company name \\
\texttt{industry\_name} & string  & Industry the company belongs to \\
\texttt{sector\_name}   & string  & Sector the company belongs to \\
\bottomrule
\end{tabular}
\end{center}

\subsubsection*{HTTP Status Codes}

\begin{center}
\small
\begin{tabular}{l p{8cm}}
\toprule
\textbf{Status} & \textbf{Description} \\
\midrule
\texttt{200 OK} & Returns array of all companies \\
\texttt{500 Internal Server Error} & Unexpected database error \\
\bottomrule
\end{tabular}
\end{center}

%───────────────────────────────────────────────────────────
\newpage
\subsection{Route 2: Search Companies}
%───────────────────────────────────────────────────────────

\route{GET}{/api/companies/search}

\medskip\noindent
\textbf{Description:} Searching companies by ticker or company name using a case-insensitive prefix match. Intended for autocomplete. Maps to Milestone~3 Query~2.

\subsubsection*{Request Parameters}

\begin{center}
\small
\begin{tabular}{l l l l p{5cm}}
\toprule
\textbf{Name} & \textbf{Param Type} & \textbf{Data Type} & \textbf{Required?} & \textbf{Description} \\
\midrule
\texttt{q} & query & string & Required & Search prefix (e.g., \texttt{AAPL} or \texttt{Apple}) \\
\bottomrule
\end{tabular}
\end{center}

\subsubsection*{SQL Query}
\begin{lstlisting}[style=sql,basicstyle=\ttfamily\scriptsize]
SELECT company_id, ticker, company_name
FROM company
WHERE ticker ILIKE ($1 || '%')
   OR company_name ILIKE ($1 || '%')
ORDER BY ticker
\end{lstlisting}

\subsubsection*{Response Schema}

\begin{center}
\small
\begin{tabular}{l l p{7cm}}
\toprule
\textbf{Field} & \textbf{Data Type} & \textbf{Description} \\
\midrule
\texttt{company\_id}   & integer & Internal company identifier \\
\texttt{ticker}        & string  & Stock ticker symbol \\
\texttt{company\_name} & string  & Full company name \\
\bottomrule
\end{tabular}
\end{center}

\subsubsection*{HTTP Status Codes}

\begin{center}
\small
\begin{tabular}{l p{8cm}}
\toprule
\textbf{Status} & \textbf{Description} \\
\midrule
\texttt{200 OK} & Returns array of matched companies (may be empty) \\
\texttt{400 Bad Request} & \texttt{q} param is missing or empty \\
\texttt{500 Internal Server Error} & Unexpected database error \\
\bottomrule
\end{tabular}
\end{center}

%───────────────────────────────────────────────────────────
\newpage
\subsection{Route 3: Company Profile}
%───────────────────────────────────────────────────────────

\route{GET}{/api/companies/:ticker}

\medskip\noindent
\textbf{Description:} Fetching a single company's profile including sector, industry, exchange, and CIK.

\subsubsection*{Request Parameters}

\begin{center}
\small
\begin{tabular}{l l l l p{5cm}}
\toprule
\textbf{Name} & \textbf{Param Type} & \textbf{Data Type} & \textbf{Required?} & \textbf{Description} \\
\midrule
\texttt{ticker} & path & string & Required & Stock ticker symbol (e.g., \texttt{AAPL}) \\
\bottomrule
\end{tabular}
\end{center}

\subsubsection*{SQL Query}
\begin{lstlisting}[style=sql,basicstyle=\ttfamily\scriptsize]
SELECT c.company_id, c.ticker, c.company_name,
       c.exchange, c.cik,
       i.industry_name, s.sector_name
FROM company c
JOIN industry i ON c.industry_id = i.industry_id
JOIN sector s   ON i.sector_id   = s.sector_id
WHERE c.ticker = $1
\end{lstlisting}

\subsubsection*{Response Schema}

\begin{center}
\small
\begin{tabular}{l l p{7cm}}
\toprule
\textbf{Field} & \textbf{Data Type} & \textbf{Description} \\
\midrule
\texttt{company\_id}    & integer & Internal company identifier \\
\texttt{ticker}         & string  & Stock ticker symbol \\
\texttt{company\_name}  & string  & Full company name \\
\texttt{exchange}       & string  & Listing exchange \\
\texttt{cik}            & string  & SEC CIK identifier \\
\texttt{industry\_name} & string  & Industry name \\
\texttt{sector\_name}   & string  & Sector name \\
\bottomrule
\end{tabular}
\end{center}

\subsubsection*{HTTP Status Codes}

\begin{center}
\small
\begin{tabular}{l p{8cm}}
\toprule
\textbf{Status} & \textbf{Description} \\
\midrule
\texttt{200 OK} & Company profile returned \\
\texttt{404 Not Found} & Ticker does not exist in the database \\
\texttt{500 Internal Server Error} & Unexpected database error \\
\bottomrule
\end{tabular}
\end{center}

%───────────────────────────────────────────────────────────
\newpage
\subsection{Route 4: Historical Prices}
%───────────────────────────────────────────────────────────

\route{GET}{/api/companies/:ticker/prices}

\medskip\noindent
\textbf{Description:} Retrieving historical OHLCV price data for a company, ordered by date. Used for rendering time-series charts. Maps to Milestone~3 Query~4.

\subsubsection*{Request Parameters}

\begin{center}
\small
\begin{tabular}{l l l l p{5cm}}
\toprule
\textbf{Name} & \textbf{Param Type} & \textbf{Data Type} & \textbf{Required?} & \textbf{Description} \\
\midrule
\texttt{ticker} & path & string & Required & Stock ticker symbol \\
\bottomrule
\end{tabular}
\end{center}

\subsubsection*{SQL Query}
\begin{lstlisting}[style=sql,basicstyle=\ttfamily\scriptsize]
SELECT c.company_name, c.ticker,
       pd.trading_date, pd.open, pd.high, pd.low,
       pd.close, pd.adj_close, pd.volume
FROM price_daily pd
JOIN company c ON c.company_id = pd.company_id
WHERE c.ticker = $1
ORDER BY pd.trading_date
\end{lstlisting}

\subsubsection*{Response Schema}

\begin{center}
\small
\begin{tabular}{l l p{7cm}}
\toprule
\textbf{Field} & \textbf{Data Type} & \textbf{Description} \\
\midrule
\texttt{company\_name}  & string  & Company name \\
\texttt{ticker}         & string  & Ticker symbol \\
\texttt{trading\_date}  & date    & Trading date \\
\texttt{open}           & float   & Opening price \\
\texttt{high}           & float   & Day high \\
\texttt{low}            & float   & Day low \\
\texttt{close}          & float   & Closing price \\
\texttt{adj\_close}     & float   & Adjusted close \\
\texttt{volume}         & integer & Trading volume \\
\bottomrule
\end{tabular}
\end{center}

\subsubsection*{HTTP Status Codes}

\begin{center}
\small
\begin{tabular}{l p{8cm}}
\toprule
\textbf{Status} & \textbf{Description} \\
\midrule
\texttt{200 OK} & Returns array of daily price records \\
\texttt{404 Not Found} & Ticker does not exist \\
\texttt{500 Internal Server Error} & Unexpected database error \\
\bottomrule
\end{tabular}
\end{center}

%───────────────────────────────────────────────────────────
\newpage
\subsection{Route 5: Top 10 Gainers (Complex)}
%───────────────────────────────────────────────────────────

\route{GET}{/api/stocks/top-gainers}

\medskip\noindent
\textbf{Description:} Returning the top 10 stocks by percentage gain on a given trading date compared to the previous day. Using CTEs and a multi-table join. Maps to Milestone~3 Query~3.

\subsubsection*{Request Parameters}

\begin{center}
\small
\begin{tabular}{l l l l p{5cm}}
\toprule
\textbf{Name} & \textbf{Param Type} & \textbf{Data Type} & \textbf{Required?} & \textbf{Description} \\
\midrule
\texttt{date} & query & date (YYYY-MM-DD) & Required & The trading date to evaluate \\
\bottomrule
\end{tabular}
\end{center}

\subsubsection*{SQL Query}
\begin{lstlisting}[style=sql,basicstyle=\ttfamily\scriptsize]
WITH latest AS (
    SELECT DISTINCT company_id, close, trading_date
    FROM price_daily
    WHERE trading_date = $1
),
prev AS (
    SELECT DISTINCT company_id, close AS prev_close
    FROM price_daily
    WHERE trading_date = $1::date - INTERVAL '1 day'
)
SELECT c.ticker, l.close, p.prev_close,
       ((l.close - p.prev_close) / p.prev_close) * 100
           AS pct_change
FROM latest l
JOIN prev p    ON l.company_id = p.company_id
JOIN company c ON c.company_id = l.company_id
ORDER BY pct_change DESC
LIMIT 10
\end{lstlisting}

\subsubsection*{Response Schema}

\begin{center}
\small
\begin{tabular}{l l p{7cm}}
\toprule
\textbf{Field} & \textbf{Data Type} & \textbf{Description} \\
\midrule
\texttt{ticker}      & string & Stock ticker symbol \\
\texttt{close}       & float  & Closing price on the selected date \\
\texttt{prev\_close} & float  & Previous day's closing price \\
\texttt{pct\_change} & float  & Percentage change from previous close \\
\bottomrule
\end{tabular}
\end{center}

\subsubsection*{HTTP Status Codes}

\begin{center}
\small
\begin{tabular}{l p{8cm}}
\toprule
\textbf{Status} & \textbf{Description} \\
\midrule
\texttt{200 OK} & Returns top 10 gainers \\
\texttt{400 Bad Request} & \texttt{date} param is missing or not a valid date \\
\texttt{500 Internal Server Error} & Unexpected database error \\
\bottomrule
\end{tabular}
\end{center}

%───────────────────────────────────────────────────────────
\newpage
\subsection{Route 6: Companies by Sector}
%───────────────────────────────────────────────────────────

\route{GET}{/api/sectors/:sector/companies}

\medskip\noindent
\textbf{Description:} Listing all companies belonging to a given sector, with their ticker and industry. Maps to Milestone~3 Query~5.

\subsubsection*{Request Parameters}

\begin{center}
\small
\begin{tabular}{l l l l p{5cm}}
\toprule
\textbf{Name} & \textbf{Param Type} & \textbf{Data Type} & \textbf{Required?} & \textbf{Description} \\
\midrule
\texttt{sector} & path & string & Required & Sector name (e.g., \texttt{Technology}) \\
\bottomrule
\end{tabular}
\end{center}

\subsubsection*{SQL Query}
\begin{lstlisting}[style=sql,basicstyle=\ttfamily\scriptsize]
SELECT c.company_id, c.ticker, c.company_name
FROM company c
JOIN industry i ON c.industry_id = i.industry_id
JOIN sector s   ON i.sector_id   = s.sector_id
WHERE s.sector_name = $1
ORDER BY c.ticker
\end{lstlisting}

\subsubsection*{Response Schema}

\begin{center}
\small
\begin{tabular}{l l p{7cm}}
\toprule
\textbf{Field} & \textbf{Data Type} & \textbf{Description} \\
\midrule
\texttt{company\_id}   & integer & Internal company identifier \\
\texttt{ticker}        & string  & Stock ticker symbol \\
\texttt{company\_name} & string  & Full company name \\
\bottomrule
\end{tabular}
\end{center}

\subsubsection*{HTTP Status Codes}

\begin{center}
\small
\begin{tabular}{l p{8cm}}
\toprule
\textbf{Status} & \textbf{Description} \\
\midrule
\texttt{200 OK} & Returns array of companies in the sector (may be empty) \\
\texttt{500 Internal Server Error} & Unexpected database error \\
\bottomrule
\end{tabular}
\end{center}

%───────────────────────────────────────────────────────────
\newpage
\subsection{Route 7: Company News}
%───────────────────────────────────────────────────────────

\route{GET}{/api/companies/:ticker/news}

\medskip\noindent
\textbf{Description:} Retrieving the 10 most recent news articles mentioning a company. Joining \texttt{news\_article} and \texttt{article\_mention}. Maps to Milestone~3 Query~6.

\subsubsection*{Request Parameters}

\begin{center}
\small
\begin{tabular}{l l l l p{5cm}}
\toprule
\textbf{Name} & \textbf{Param Type} & \textbf{Data Type} & \textbf{Required?} & \textbf{Description} \\
\midrule
\texttt{ticker} & path & string & Required & Stock ticker symbol \\
\bottomrule
\end{tabular}
\end{center}

\subsubsection*{SQL Query}
\begin{lstlisting}[style=sql,basicstyle=\ttfamily\scriptsize]
SELECT c.company_name, c.ticker,
       na.title, na.source, na.published_at, na.url
FROM news_article na
JOIN article_mention am ON na.article_id = am.article_id
JOIN company c          ON c.company_id  = am.company_id
WHERE c.ticker = $1
ORDER BY na.published_at DESC
LIMIT 10
\end{lstlisting}

\subsubsection*{Response Schema}

\begin{center}
\small
\begin{tabular}{l l p{7cm}}
\toprule
\textbf{Field} & \textbf{Data Type} & \textbf{Description} \\
\midrule
\texttt{company\_name} & string      & Company name \\
\texttt{ticker}        & string      & Ticker symbol \\
\texttt{title}         & string      & Article headline \\
\texttt{source}        & string      & News source name \\
\texttt{published\_at} & timestamptz & Publication timestamp \\
\texttt{url}           & string      & Article URL \\
\bottomrule
\end{tabular}
\end{center}

\subsubsection*{HTTP Status Codes}

\begin{center}
\small
\begin{tabular}{l p{8cm}}
\toprule
\textbf{Status} & \textbf{Description} \\
\midrule
\texttt{200 OK} & Returns array of news articles (may be empty) \\
\texttt{404 Not Found} & Ticker does not exist \\
\texttt{500 Internal Server Error} & Unexpected database error \\
\bottomrule
\end{tabular}
\end{center}

%───────────────────────────────────────────────────────────
\newpage
\subsection{Route 8: Trending Companies in News (Complex)}
%───────────────────────────────────────────────────────────

\route{GET}{/api/news/trending}

\medskip\noindent
\textbf{Description:} Returning the top 10 most-mentioned companies in news articles over the past 30 days. Using aggregation with \texttt{GROUP BY} and a date-range filter. Maps to Milestone~3 Query~7.

\subsubsection*{Request Parameters}

None.

\subsubsection*{SQL Query}
\begin{lstlisting}[style=sql,basicstyle=\ttfamily\scriptsize]
SELECT c.ticker, COUNT(*) AS mention_count
FROM article_mention am
JOIN company c      ON am.company_id = c.company_id
JOIN news_article na ON na.article_id = am.article_id
WHERE na.published_at >= CURRENT_DATE - INTERVAL '30 days'
GROUP BY c.ticker
ORDER BY mention_count DESC
LIMIT 10
\end{lstlisting}

\subsubsection*{Response Schema}

\begin{center}
\small
\begin{tabular}{l l p{7cm}}
\toprule
\textbf{Field} & \textbf{Data Type} & \textbf{Description} \\
\midrule
\texttt{ticker}         & string  & Stock ticker symbol \\
\texttt{mention\_count} & integer & Number of article mentions in the past 30 days \\
\bottomrule
\end{tabular}
\end{center}

\subsubsection*{HTTP Status Codes}

\begin{center}
\small
\begin{tabular}{l p{8cm}}
\toprule
\textbf{Status} & \textbf{Description} \\
\midrule
\texttt{200 OK} & Returns top 10 trending companies \\
\texttt{500 Internal Server Error} & Unexpected database error \\
\bottomrule
\end{tabular}
\end{center}

%───────────────────────────────────────────────────────────
\newpage
\subsection{Route 9: Sector Returns (Complex)}
%───────────────────────────────────────────────────────────

\route{GET}{/api/sectors/returns}

\medskip\noindent
\textbf{Description:} Computing the average daily return by sector on a given date compared to the previous day. Using CTEs, multi-table joins, and aggregation across sector/industry/company/price tables. Maps to Milestone~3 Query~8.

\subsubsection*{Request Parameters}

\begin{center}
\small
\begin{tabular}{l l l l p{5cm}}
\toprule
\textbf{Name} & \textbf{Param Type} & \textbf{Data Type} & \textbf{Required?} & \textbf{Description} \\
\midrule
\texttt{date}   & query & date (YYYY-MM-DD) & Required & The trading date to evaluate \\
\texttt{sector} & query & string            & Optional & Filter to a specific sector \\
\bottomrule
\end{tabular}
\end{center}

\subsubsection*{SQL Query}
\begin{lstlisting}[style=sql,basicstyle=\ttfamily\scriptsize]
WITH latest AS (
    SELECT DISTINCT company_id, close
    FROM price_daily
    WHERE trading_date = $1
),
prev AS (
    SELECT DISTINCT company_id, close AS prev_close
    FROM price_daily
    WHERE trading_date = $1::date - INTERVAL '1 day'
)
SELECT s.sector_name,
       AVG((l.close - p.prev_close) / p.prev_close * 100)
           AS avg_return_pct
FROM latest l
JOIN prev p    ON l.company_id = p.company_id
JOIN company c ON c.company_id = l.company_id
JOIN industry i ON c.industry_id = i.industry_id
JOIN sector s   ON i.sector_id   = s.sector_id
WHERE ($2 IS NULL OR s.sector_name = $2)
GROUP BY s.sector_name
ORDER BY avg_return_pct DESC
\end{lstlisting}

\subsubsection*{Response Schema}

\begin{center}
\small
\begin{tabular}{l l p{7cm}}
\toprule
\textbf{Field} & \textbf{Data Type} & \textbf{Description} \\
\midrule
\texttt{sector\_name}     & string & Sector name \\
\texttt{avg\_return\_pct} & float  & Average return percentage across companies in the sector \\
\bottomrule
\end{tabular}
\end{center}

\subsubsection*{HTTP Status Codes}

\begin{center}
\small
\begin{tabular}{l p{8cm}}
\toprule
\textbf{Status} & \textbf{Description} \\
\midrule
\texttt{200 OK} & Returns sector return data \\
\texttt{400 Bad Request} & \texttt{date} param is missing or not a valid date \\
\texttt{500 Internal Server Error} & Unexpected database error \\
\bottomrule
\end{tabular}
\end{center}

%───────────────────────────────────────────────────────────
\newpage
\subsection{Route 10: List Sectors}
%───────────────────────────────────────────────────────────

\route{GET}{/api/sectors}

\medskip\noindent
\textbf{Description:} Returning all sectors. Used for populating filter dropdowns in the frontend.

\subsubsection*{Request Parameters}

None.

\subsubsection*{SQL Query}
\begin{lstlisting}[style=sql,basicstyle=\ttfamily\scriptsize]
SELECT sector_id, sector_name
FROM sector
ORDER BY sector_name
\end{lstlisting}

\subsubsection*{Response Schema}

\begin{center}
\small
\begin{tabular}{l l p{7cm}}
\toprule
\textbf{Field} & \textbf{Data Type} & \textbf{Description} \\
\midrule
\texttt{sector\_id}   & integer & Internal sector identifier \\
\texttt{sector\_name} & string  & Sector name \\
\bottomrule
\end{tabular}
\end{center}

\subsubsection*{HTTP Status Codes}

\begin{center}
\small
\begin{tabular}{l p{8cm}}
\toprule
\textbf{Status} & \textbf{Description} \\
\midrule
\texttt{200 OK} & Returns array of all sectors \\
\texttt{500 Internal Server Error} & Unexpected database error \\
\bottomrule
\end{tabular}
\end{center}

%═══════════════════════════════════════════════════════════
\newpage
\section{Auxiliary Routes}
%═══════════════════════════════════════════════════════════

%───────────────────────────────────────────────────────────
\subsection{Route A1: Create User}
%───────────────────────────────────────────────────────────

\route{POST}{/api/users}

\medskip\noindent
\textbf{Description:} Creating a new user account. Supporting future watchlist and preference features.

\subsubsection*{Request Parameters}

\begin{center}
\small
\begin{tabular}{l l l l p{5cm}}
\toprule
\textbf{Name} & \textbf{Param Type} & \textbf{Data Type} & \textbf{Required?} & \textbf{Description} \\
\midrule
\texttt{username} & body & string & Required & Desired username \\
\texttt{email}    & body & string & Required & User email address \\
\texttt{password} & body & string & Required & Password (hashed server-side) \\
\bottomrule
\end{tabular}
\end{center}

\subsubsection*{SQL Query}
\begin{lstlisting}[style=sql,basicstyle=\ttfamily\scriptsize]
INSERT INTO users (username, email, password_hash)
VALUES ($1, $2, $3)
RETURNING user_id, username
\end{lstlisting}

\subsubsection*{Response Schema}

\begin{center}
\small
\begin{tabular}{l l p{7cm}}
\toprule
\textbf{Field} & \textbf{Data Type} & \textbf{Description} \\
\midrule
\texttt{user\_id}  & integer & Newly created user's ID \\
\texttt{username}  & string  & Confirmed username \\
\bottomrule
\end{tabular}
\end{center}

\subsubsection*{HTTP Status Codes}

\begin{center}
\small
\begin{tabular}{l p{8cm}}
\toprule
\textbf{Status} & \textbf{Description} \\
\midrule
\texttt{201 Created} & User account created \\
\texttt{400 Bad Request} & Required field missing \\
\texttt{409 Conflict} & Username or email already exists \\
\texttt{500 Internal Server Error} & Unexpected database error \\
\bottomrule
\end{tabular}
\end{center}

%───────────────────────────────────────────────────────────
\newpage
\subsection{Route A2: User Login}
%───────────────────────────────────────────────────────────

\route{POST}{/api/users/login}

\medskip\noindent
\textbf{Description:} Authenticating an existing user by username and password.

\subsubsection*{Request Parameters}

\begin{center}
\small
\begin{tabular}{l l l l p{5cm}}
\toprule
\textbf{Name} & \textbf{Param Type} & \textbf{Data Type} & \textbf{Required?} & \textbf{Description} \\
\midrule
\texttt{username} & body & string & Required & Existing username \\
\texttt{password} & body & string & Required & Password to verify \\
\bottomrule
\end{tabular}
\end{center}

\subsubsection*{SQL Query}
\begin{lstlisting}[style=sql,basicstyle=\ttfamily\scriptsize]
SELECT user_id, password_hash
FROM users
WHERE username = $1
\end{lstlisting}

\subsubsection*{Response Schema}

\begin{center}
\small
\begin{tabular}{l l p{7cm}}
\toprule
\textbf{Field} & \textbf{Data Type} & \textbf{Description} \\
\midrule
\texttt{token}    & string  & Session token \\
\texttt{user\_id} & integer & Authenticated user's ID \\
\bottomrule
\end{tabular}
\end{center}

\subsubsection*{HTTP Status Codes}

\begin{center}
\small
\begin{tabular}{l p{8cm}}
\toprule
\textbf{Status} & \textbf{Description} \\
\midrule
\texttt{200 OK} & Login successful \\
\texttt{400 Bad Request} & Required field missing \\
\texttt{401 Unauthorized} & Incorrect password \\
\texttt{404 Not Found} & Username does not exist \\
\texttt{500 Internal Server Error} & Unexpected database error \\
\bottomrule
\end{tabular}
\end{center}

\end{document}
